This appendix records the coordinate checks used for the basic hyperbolic
arithmetic expression space.  Throughout, \(\mu,\lambda>0\),
\[
 G=\{(x,y):x\in\mathbb R,\ y>0\},
 \qquad
 (x,y)(x',y')=(x+yx',yy'),
\]
and
\[
 g_{\mu,\lambda}
 =\frac1{y^2}
  \left(\frac{dx^2}{\mu^2}+\frac{dy^2}{\lambda^2}\right),
 \qquad
 a=-\frac xy.
\]

\subsection{Group coordinates and the invariant frame}

The identity of \(G\) is \((0,1)\), and
\[
 (x,y)^{-1}=\left(-\frac xy,\frac1y\right).
\]
For \(g_0=(x_0,y_0)\), left translation and its differential are
\[
 L_{g_0}(x,y)=(x_0+y_0x,y_0y),
 \qquad
 dL_{g_0}(\partial_x)=y_0\partial_x,
 \qquad
 dL_{g_0}(\partial_y)=y_0\partial_y.
\]
It follows that \(y\partial_x\) and \(y\partial_y\) are left-invariant.
The normalized arithmetic frame
\[
 X_u=-\mu y\partial_x,
 \qquad
 X_v=-\lambda y\partial_y
\]
is therefore left-invariant, and its dual coframe is
\[
 -\frac{dx}{\mu y},
 \qquad
 -\frac{dy}{\lambda y}.
\]
The sum of their squares is \(g_{\mu,\lambda}\).  The bracket is
\begin{align*}
 [X_u,X_v]
 &=[-\mu y\partial_x,-\lambda y\partial_y]\\
 &=-\mu\lambda y\partial_x
 =\lambda X_u,
\end{align*}
and hence \([X_u,X_v]a=\mu\lambda\), as required by the general
arithmetic-frame identity.

The same invariance can be checked directly:
\begin{align*}
 L_{(x_0,y_0)}^*g_{\mu,\lambda}
 &=
 \frac1{(y_0y)^2}
 \left(
  \frac{(y_0dx)^2}{\mu^2}
  +\frac{(y_0dy)^2}{\lambda^2}
 \right)\\
 &=g_{\mu,\lambda}.
\end{align*}

For a horocyclic chart, set
\[
 u=x,\qquad y=e^{\lambda v}.
\]
Then
\[
 (u,v)(u',v')
 =\bigl(u+e^{\lambda v}u',v+v'\bigr),
\]
and
\[
 X_u=-\mu e^{\lambda v}\partial_u,
 \qquad
 X_v=-\partial_v,
\]
while
\[
 g_{\mu,\lambda}
 =e^{-2\lambda v}\frac{du^2}{\mu^2}+dv^2,
 \qquad
 a=-u e^{-\lambda v}.
\]
The letters \(u,v\) in this chart are genuine coordinates, but the arithmetic
frame is not the coordinate frame.

\subsection{Inverse metric, gradient, and the eikonal identity}

In the \((x,y)\)-chart,
\[
 (g_{ij})
 =
 \begin{pmatrix}
  \dfrac1{\mu^2y^2}&0\\[5pt]
  0&\dfrac1{\lambda^2y^2}
 \end{pmatrix},
 \qquad
 (g^{ij})
 =
 \begin{pmatrix}
  \mu^2y^2&0\\
  0&\lambda^2y^2
 \end{pmatrix},
\]
and
\[
 \det(g_{ij})=\frac1{\mu^2\lambda^2y^4},
 \qquad
 d\operatorname{vol}_g
 =\frac{dx\wedge dy}{\mu\lambda y^2}.
\]
Since
\[
 a_x=-\frac1y,
 \qquad
 a_y=\frac{x}{y^2}=-\frac ay,
\]
the gradient is
\begin{align*}
 \nabla a
 &=g^{xx}a_x\partial_x+g^{yy}a_y\partial_y\\
 &=-\mu^2y\,\partial_x+\lambda^2x\,\partial_y.
\end{align*}
Its squared norm is
\begin{align*}
 |\nabla a|_g^2
 &=g^{xx}a_x^2+g^{yy}a_y^2\\
 &=\mu^2+\lambda^2\frac{x^2}{y^2}\\
 &=\mu^2+\lambda^2a^2.
\end{align*}
Alternatively,
\[
 X_u a=(-\mu y)\left(-\frac1y\right)=\mu,
 \qquad
 X_v a=(-\lambda y)\left(\frac{x}{y^2}\right)=\lambda a,
\]
so the framed and eikonal verifications agree.

\subsection{Curvature and completeness}

Two calculations give the curvature normalization.  First set
\[
 X=\frac{\lambda}{\mu}x.
\]
This is a global diffeomorphism of upper half-planes, and
\[
 g_{\mu,\lambda}
 =\frac1{\lambda^2}
  \frac{dX^2+dy^2}{y^2}.
\]
The standard metric \((dX^2+dy^2)/y^2\) is complete with curvature
\(-1\).  Constant rescaling by \(\lambda^{-2}\) gives a complete metric
of curvature \(-\lambda^2\).

For a direct Cartan calculation, use the positively oriented orthonormal coframe
\[
 \omega^1=\frac{dx}{\mu y},
 \qquad
 \omega^2=\frac{dy}{\lambda y}.
\]
Then
\[
 d\omega^1=\lambda\,\omega^1\wedge\omega^2,
 \qquad
 d\omega^2=0.
\]
The Levi--Civita connection form
\(\omega^1_{\ 2}=-\lambda\omega^1\) satisfies
\[
 d\omega^1=-\omega^1_{\ 2}\wedge\omega^2,
 \qquad
 d\omega^2=\omega^1_{\ 2}\wedge\omega^1.
\]
Its curvature form is
\[
 d\omega^1_{\ 2}
 =-\lambda\,d\omega^1
 =-\lambda^2\,\omega^1\wedge\omega^2.
\]
Thus \(K=-\lambda^2\), consistently with the coordinate rescaling.

\subsection{The Laplace--Beltrami operator}

With the divergence-of-gradient convention,
\[
 \Delta_g f
 =\frac1{\sqrt{\det g}}\,
  \partial_i\!\left(\sqrt{\det g}\,g^{ij}\partial_jf\right).
\]
Here
\[
 \sqrt{\det g}=\frac1{\mu\lambda y^2},
 \quad
 \sqrt{\det g}\,g^{xx}=\frac{\mu}{\lambda},
 \quad
 \sqrt{\det g}\,g^{yy}=\frac{\lambda}{\mu}.
\]
The last two coefficients are constant, so
\begin{equation}\label{eq:appendix-hyperbolic-laplacian}
 \Delta_{g_{\mu,\lambda}}f
 =\mu^2y^2f_{xx}+\lambda^2y^2f_{yy}.
\end{equation}
For \(a=-x/y\),
\[
 a_{xx}=0,
 \qquad
 a_{yy}=-\frac{2x}{y^3}=\frac{2a}{y^2}.
\]
Equation~\eqref{eq:appendix-hyperbolic-laplacian} therefore gives
\[
 \Delta_{g_{\mu,\lambda}}a=2\lambda^2a.
\]

\subsection{Grid maps and their metric pullbacks}

For
\[
 \mathsf X_s(x,y)=(x-sy,y),
 \qquad
 \mathsf Y_k(x,y)=\left(x,\frac yk\right),
\]
one computes
\[
 (a\circ\mathsf X_s)(x,y)
 =-\frac{x-sy}{y}=a(x,y)+s
\]
and
\[
 (a\circ\mathsf Y_k)(x,y)
 =-\frac{x}{y/k}=k\,a(x,y).
\]
On differentials,
\[
 \mathsf X_s^*dx=dx-s\,dy,\qquad
 \mathsf X_s^*dy=dy,
\]
and the transformed \(y\)-coordinate is still \(y\).  Hence
\[
 \mathsf X_s^*g_{\mu,\lambda}
 =\frac1{y^2}
  \left(
   \frac{(dx-s\,dy)^2}{\mu^2}
   +\frac{dy^2}{\lambda^2}
  \right).
\]
For \(\mathsf Y_k\), the transformed coordinate and differentials are
\[
 y'=\frac yk,\qquad dx'=dx,\qquad dy'=\frac{dy}{k},
\]
so
\[
 \mathsf Y_k^*g_{\mu,\lambda}
 =\frac1{y^2}
  \left(
   \frac{k^2dx^2}{\mu^2}
   +\frac{dy^2}{\lambda^2}
  \right).
\]
The cross term in the first formula and the unequal horizontal coefficient in
the second show that these grid maps are not general isometries.

\subsection{The Baumslag--Solitar composition check}

For a positive integer \(n\),
\[
 \mathsf X_s^n(x,y)=(x-nsy,y).
\]
Using the standard right-to-left convention for composition,
\begin{align*}
 (\mathsf Y_n^{-1}\circ\mathsf X_s^n\circ\mathsf Y_n)(x,y)
 &=(\mathsf Y_n^{-1}\circ\mathsf X_s^n)
   \left(x,\frac yn\right)\\
 &=\mathsf Y_n^{-1}(x-sy,y/n)\\
 &=(x-sy,y)\\
 &=\mathsf X_s(x,y).
\end{align*}
Equivalently,
\[
 \mathsf Y_n\circ\mathsf X_s\circ\mathsf Y_n^{-1}
 =\mathsf X_s^n.
\]
These are identities in the diffeomorphism group of \(\mathbb H^2\).  The metric
pullback calculation above is why they are stated as grid relations rather than
as relations in the hyperbolic isometry group.
